\chapter{Unit 1: Introduction to Environment and Sustainable Development}

\section{Section A: Multiple Choice Practice Questions (60 Curated Problems)}

\noindent\textbf{Q1. What does the term environment refer to?} \hfill [Tag: Introduction to environment | Easy]
\begin{enumerate}[label=(\alph*),topsep=2pt,itemsep=1pt,partopsep=0pt,parsep=0pt]
    \item {Only the weather of a place}
    \item {The surroundings in which organisms live}
    \item {Only the plants in a forest}
    \item {Only the buildings around people}
\end{enumerate}
\vspace{2mm}

\noindent\textbf{Q2. Which of the following is a natural component of the environment?} \hfill [Tag: Introduction to environment | Easy]
\begin{enumerate}[label=(\alph*),topsep=2pt,itemsep=1pt,partopsep=0pt,parsep=0pt]
    \item {Air}
    \item {Bridge}
    \item {Factory}
    \item {Road}
\end{enumerate}
\vspace{2mm}

\noindent\textbf{Q3. Which component of the environment includes plants, animals, and microorganisms?} \hfill [Tag: Components of environment | Easy]
\begin{enumerate}[label=(\alph*),topsep=2pt,itemsep=1pt,partopsep=0pt,parsep=0pt]
    \item {Abiotic component}
    \item {Geological component}
    \item {Physical component}
    \item {Biotic component}
\end{enumerate}
\vspace{2mm}

\noindent\textbf{Q4. Which of the following is an abiotic component of the environment?} \hfill [Tag: Components of environment | Easy]
\begin{enumerate}[label=(\alph*),topsep=2pt,itemsep=1pt,partopsep=0pt,parsep=0pt]
    \item {Trees}
    \item {Birds}
    \item {Water}
    \item {Bacteria}
\end{enumerate}
\vspace{2mm}

\noindent\textbf{Q5. Which sphere of Earth consists mainly of air?} \hfill [Tag: Spheres of earth | Easy]
\begin{enumerate}[label=(\alph*),topsep=2pt,itemsep=1pt,partopsep=0pt,parsep=0pt]
    \item {Atmosphere}
    \item {Hydrosphere}
    \item {Biosphere}
    \item {Lithosphere}
\end{enumerate}
\vspace{2mm}

\noindent\textbf{Q6. The hydrosphere includes which of the following?} \hfill [Tag: Spheres of earth | Easy]
\begin{enumerate}[label=(\alph*),topsep=2pt,itemsep=1pt,partopsep=0pt,parsep=0pt]
    \item {All rocks on Earth}
    \item {All living organisms}
    \item {All water on Earth}
    \item {All gases around Earth}
\end{enumerate}
\vspace{2mm}

\noindent\textbf{Q7. Which sphere includes the solid outer part of Earth?} \hfill [Tag: Spheres of earth | Easy]
\begin{enumerate}[label=(\alph*),topsep=2pt,itemsep=1pt,partopsep=0pt,parsep=0pt]
    \item {Lithosphere}
    \item {Atmosphere}
    \item {Hydrosphere}
    \item {Biosphere}
\end{enumerate}
\vspace{2mm}

\noindent\textbf{Q8. The biosphere is the part of Earth where:} \hfill [Tag: Spheres of earth | Easy]
\begin{enumerate}[label=(\alph*),topsep=2pt,itemsep=1pt,partopsep=0pt,parsep=0pt]
    \item {Life exists}
    \item {Only rocks occur}
    \item {Only clouds are formed}
    \item {Only water is stored}
\end{enumerate}
\vspace{2mm}

\noindent\textbf{Q9. Why is environmental studies called a multidisciplinary subject?} \hfill [Tag: Multidisciplinary nature, scope and importance of environmental studies | Easy]
\begin{enumerate}[label=(\alph*),topsep=2pt,itemsep=1pt,partopsep=0pt,parsep=0pt]
    \item {It excludes social and economic issues}
    \item {It studies only one type of organism}
    \item {It uses knowledge from many subjects}
    \item {It focuses only on chemical reactions}
\end{enumerate}
\vspace{2mm}

\noindent\textbf{Q10. Which field helps environmental studies understand human population and communities?} \hfill [Tag: Multidisciplinary nature, scope and importance of environmental studies | Easy]
\begin{enumerate}[label=(\alph*),topsep=2pt,itemsep=1pt,partopsep=0pt,parsep=0pt]
    \item {Sociology}
    \item {Trigonometry}
    \item {Linguistics}
    \item {Astronomy}
\end{enumerate}
\vspace{2mm}

\noindent\textbf{Q11. What is one important aim of environmental studies?} \hfill [Tag: Multidisciplinary nature, scope and importance of environmental studies | Easy]
\begin{enumerate}[label=(\alph*),topsep=2pt,itemsep=1pt,partopsep=0pt,parsep=0pt]
    \item {To increase resource wastage}
    \item {To discourage public awareness}
    \item {To promote environmental protection}
    \item {To remove all natural habitats}
\end{enumerate}
\vspace{2mm}

\noindent\textbf{Q12. Environmental studies is important because it helps people:} \hfill [Tag: Multidisciplinary nature, scope and importance of environmental studies | Easy]
\begin{enumerate}[label=(\alph*),topsep=2pt,itemsep=1pt,partopsep=0pt,parsep=0pt]
    \item {Understand environmental problems}
    \item {Use resources without limits}
    \item {Ignore pollution problems}
    \item {Avoid learning about ecosystems}
\end{enumerate}
\vspace{2mm}

\noindent\textbf{Q13. What is the basic meaning of sustainability?} \hfill [Tag: Concept of sustainability | Easy]
\begin{enumerate}[label=(\alph*),topsep=2pt,itemsep=1pt,partopsep=0pt,parsep=0pt]
    \item {Avoiding every form of development}
    \item {Increasing waste in natural systems}
    \item {Using resources responsibly over time}
    \item {Using all resources immediately}
\end{enumerate}
\vspace{2mm}

\noindent\textbf{Q14. Sustainable development aims to meet present needs without:} \hfill [Tag: Sustainable development | Easy]
\begin{enumerate}[label=(\alph*),topsep=2pt,itemsep=1pt,partopsep=0pt,parsep=0pt]
    \item {Harming future generations' needs}
    \item {Improving present living conditions}
    \item {Using any natural resources}
    \item {Supporting economic activities}
\end{enumerate}
\vspace{2mm}

\noindent\textbf{Q15. Which action supports sustainable development?} \hfill [Tag: Sustainable development | Easy]
\begin{enumerate}[label=(\alph*),topsep=2pt,itemsep=1pt,partopsep=0pt,parsep=0pt]
    \item {Dumping waste into rivers}
    \item {Using renewable energy}
    \item {Wasting clean drinking water}
    \item {Cutting forests without replanting}
\end{enumerate}
\vspace{2mm}

\noindent\textbf{Q16. What is carrying capacity?} \hfill [Tag: Carrying capacity | Easy]
\begin{enumerate}[label=(\alph*),topsep=2pt,itemsep=1pt,partopsep=0pt,parsep=0pt]
    \item {The speed of population growth}
    \item {The amount of rainfall in a region}
    \item {The total number of species on Earth}
    \item {The largest population an environment can support}
\end{enumerate}
\vspace{2mm}

\noindent\textbf{Q17. What may happen when a population exceeds its carrying capacity?} \hfill [Tag: Carrying capacity | Easy]
\begin{enumerate}[label=(\alph*),topsep=2pt,itemsep=1pt,partopsep=0pt,parsep=0pt]
    \item {Competition between organisms ends}
    \item {Resources automatically increase}
    \item {Pollution always disappears}
    \item {Resources may become depleted}
\end{enumerate}
\vspace{2mm}

\noindent\textbf{Q18. Which of the following is commonly considered a pillar of sustainability?} \hfill [Tag: Pillars of sustainability | Easy]
\begin{enumerate}[label=(\alph*),topsep=2pt,itemsep=1pt,partopsep=0pt,parsep=0pt]
    \item {Entertainment}
    \item {Fashion}
    \item {Advertising}
    \item {Environment}
\end{enumerate}
\vspace{2mm}

\noindent\textbf{Q19. Which set correctly lists the three main pillars of sustainability?} \hfill [Tag: Pillars of sustainability | Easy]
\begin{enumerate}[label=(\alph*),topsep=2pt,itemsep=1pt,partopsep=0pt,parsep=0pt]
    \item {Energy, transport, and tourism}
    \item {Environment, economy, and society}
    \item {Science, language, and history}
    \item {Agriculture, industry, and trade}
\end{enumerate}
\vspace{2mm}

\noindent\textbf{Q20. How many Sustainable Development Goals were adopted by the United Nations?} \hfill [Tag: Sustainable development goals | Easy]
\begin{enumerate}[label=(\alph*),topsep=2pt,itemsep=1pt,partopsep=0pt,parsep=0pt]
    \item {10}
    \item {17}
    \item {twenty}
    \item {15}
\end{enumerate}
\vspace{2mm}

\noindent\textbf{Q21. A new highway reduces travel time but fragments a wildlife habitat and increases nearby noise levels. Which interpretation best reflects the environmental perspective?} \hfill [Tag: Introduction to environment | Medium]
\begin{enumerate}[label=(\alph*),topsep=2pt,itemsep=1pt,partopsep=0pt,parsep=0pt]
    \item {The project affects only the region's biological surroundings}
    \item {The project is unrelated to environment because it is economic}
    \item {The project is environmental only if air pollution increases}
    \item {The project affects interacting natural and human surroundings}
\end{enumerate}
\vspace{2mm}

\noindent\textbf{Q22. Residents report that a lake has become unsuitable for recreation after untreated sewage caused algal growth and foul odors. Which relationship is most directly illustrated?} \hfill [Tag: Introduction to environment | Medium]
\begin{enumerate}[label=(\alph*),topsep=2pt,itemsep=1pt,partopsep=0pt,parsep=0pt]
    \item {Geological activity increasing social diversity}
    \item {Climate variation controlling urban employment}
    \item {Human activity altering environmental quality}
    \item {Wildlife migration changing industrial production}
\end{enumerate}
\vspace{2mm}

\noindent\textbf{Q23. A city replaces a wetland with buildings, after which flooding increases and bird populations decline. Which environmental components are most directly involved?} \hfill [Tag: Components of environment | Medium]
\begin{enumerate}[label=(\alph*),topsep=2pt,itemsep=1pt,partopsep=0pt,parsep=0pt]
    \item {Biotic, abiotic, and cultural components}
    \item {Abiotic and cultural components only}
    \item {Abiotic and extraterrestrial components}
    \item {Biotic and cultural components only}
\end{enumerate}
\vspace{2mm}

\noindent\textbf{Q24. Which observation best demonstrates an interaction between biotic and abiotic components of the environment?} \hfill [Tag: Components of environment | Medium]
\begin{enumerate}[label=(\alph*),topsep=2pt,itemsep=1pt,partopsep=0pt,parsep=0pt]
    \item {Wind transports sand across an empty desert}
    \item {A city expands its public transport network}
    \item {Two bird species compete for nesting space}
    \item {Tree roots reduce soil erosion on a slope}
\end{enumerate}
\vspace{2mm}

\noindent\textbf{Q25. Sulfur dioxide released into the air forms acid rain that lowers the pH of a lake and harms fish. Which sequence of Earth-sphere interactions is most accurate?} \hfill [Tag: Spheres of earth | Medium]
\begin{enumerate}[label=(\alph*),topsep=2pt,itemsep=1pt,partopsep=0pt,parsep=0pt]
    \item {Biosphere to lithosphere to atmosphere}
    \item {Atmosphere to hydrosphere to biosphere}
    \item {Hydrosphere to atmosphere to lithosphere}
    \item {Lithosphere to biosphere to hydrosphere}
\end{enumerate}
\vspace{2mm}

\noindent\textbf{Q26. During heavy rain, exposed soil from a construction site enters a river and reduces aquatic plant growth. Which three spheres interact most directly?} \hfill [Tag: Spheres of earth | Medium]
\begin{enumerate}[label=(\alph*),topsep=2pt,itemsep=1pt,partopsep=0pt,parsep=0pt]
    \item {Atmosphere, lithosphere, and magnetosphere}
    \item {Biosphere, cryosphere, and magnetosphere}
    \item {Hydrosphere, atmosphere, and cryosphere}
    \item {Lithosphere, hydrosphere, and biosphere}
\end{enumerate}
\vspace{2mm}

\noindent\textbf{Q27. A town wants to control air pollution without unfairly burdening low-income workers. Which team would provide the most appropriate multidisciplinary assessment?} \hfill [Tag: Multidisciplinary nature, scope and importance of environmental studies | Medium]
\begin{enumerate}[label=(\alph*),topsep=2pt,itemsep=1pt,partopsep=0pt,parsep=0pt]
    \item {An ecologist, an economist, and a sociologist}
    \item {A botanist, a zoologist, and a microbiologist}
    \item {An architect, a historian, and a linguist}
    \item {A geologist, a physicist, and an astronomer}
\end{enumerate}
\vspace{2mm}

\noindent\textbf{Q28. Before approving a dam, planners study fish habitats, displaced communities, electricity demand, and legal requirements. Why is this approach appropriate?} \hfill [Tag: Multidisciplinary nature, scope and importance of environmental studies | Medium]
\begin{enumerate}[label=(\alph*),topsep=2pt,itemsep=1pt,partopsep=0pt,parsep=0pt]
    \item {Environmental decisions involve ecological, social, economic, and legal factors}
    \item {Environmental decisions depend mainly on biological classification and naming}
    \item {Environmental decisions become reliable only after construction is completed}
    \item {Environmental decisions should exclude effects that cannot be priced}
\end{enumerate}
\vspace{2mm}

\noindent\textbf{Q29. A forest produces about 10 , 000 tonnes of new timber each year. Which annual harvesting plan is most consistent with ecological sustainability, assuming other conditions remain stable?} \hfill [Tag: Concept of sustainability | Medium]
\begin{enumerate}[label=(\alph*),topsep=2pt,itemsep=1pt,partopsep=0pt,parsep=0pt]
    \item {Harvest 12 ,}
    \item {Harvest 8 ,}
    \item {Harvest 10 ,}
    \item {Harvest 15 ,}
\end{enumerate}
\vspace{2mm}

\noindent\textbf{Q30. A university wants to make its cafeteria more sustainable. Which action best applies a life-cycle perspective?} \hfill [Tag: Concept of sustainability | Medium]
\begin{enumerate}[label=(\alph*),topsep=2pt,itemsep=1pt,partopsep=0pt,parsep=0pt]
    \item {Measuring only the electricity used during lunch hours}
    \item {Comparing food production, transport, use, and disposal impacts}
    \item {Counting only the waste placed in cafeteria bins}
    \item {Choosing meals solely according to their purchase prices}
\end{enumerate}
\vspace{2mm}

\noindent\textbf{Q31. Which urban policy best reflects sustainable development?} \hfill [Tag: Sustainable development | Medium]
\begin{enumerate}[label=(\alph*),topsep=2pt,itemsep=1pt,partopsep=0pt,parsep=0pt]
    \item {Increasing road capacity while removing urban green spaces}
    \item {Expanding affordable transit while reducing vehicle emissions}
    \item {Restricting housing construction while promoting private vehicles}
    \item {Closing factories without providing alternative local employment}
\end{enumerate}
\vspace{2mm}

\noindent\textbf{Q32. A coastal community must choose a development proposal. Which proposal best meets present needs without undermining future generations?} \hfill [Tag: Sustainable development | Medium]
\begin{enumerate}[label=(\alph*),topsep=2pt,itemsep=1pt,partopsep=0pt,parsep=0pt]
    \item {Exclusive resorts with permanent removal of mangrove forests}
    \item {Large-scale mining with compensation after resource depletion}
    \item {Managed ecotourism with habitat limits and local employment}
    \item {Intensive tourism with unrestricted groundwater extraction}
\end{enumerate}
\vspace{2mm}

\noindent\textbf{Q33. An island has 3 , 600 , 000 liters of renewable freshwater available annually. If each resident requires 120 , 000 liters per year, what is the water-based carrying capacity?} \hfill [Tag: Carrying capacity | Medium]
\begin{enumerate}[label=(\alph*),topsep=2pt,itemsep=1pt,partopsep=0pt,parsep=0pt]
    \item {40}
    \item {20}
    \item {30}
    \item {50}
\end{enumerate}
\vspace{2mm}

\noindent\textbf{Q34. A deer population remains above the carrying capacity of a grassland for several years. What outcome is most likely if no limiting factor changes?} \hfill [Tag: Carrying capacity | Medium]
\begin{enumerate}[label=(\alph*),topsep=2pt,itemsep=1pt,partopsep=0pt,parsep=0pt]
    \item {Unlimited growth followed by habitat expansion}
    \item {Permanent equilibrium above available resource limits}
    \item {Food depletion followed by population decline}
    \item {Stable growth with steadily increasing food supply}
\end{enumerate}
\vspace{2mm}

\noindent\textbf{Q35. A farming region adopts drip irrigation and soil-moisture monitoring, allowing the same water supply to support more crops. How does this affect carrying capacity?} \hfill [Tag: Carrying capacity | Medium]
\begin{enumerate}[label=(\alph*),topsep=2pt,itemsep=1pt,partopsep=0pt,parsep=0pt]
    \item {It disappears because water is no longer a limiting factor}
    \item {It decreases because technology always raises consumption}
    \item {It remains fixed because carrying capacity cannot change}
    \item {It increases because water is used more efficiently}
\end{enumerate}
\vspace{2mm}

\noindent\textbf{Q36. A solar project lowers emissions and earns a profit but displaces residents without consultation or compensation. Which sustainability pillar is most clearly neglected?} \hfill [Tag: Pillars of sustainability | Medium]
\begin{enumerate}[label=(\alph*),topsep=2pt,itemsep=1pt,partopsep=0pt,parsep=0pt]
    \item {Social sustainability}
    \item {Environmental sustainability}
    \item {Technological sustainability}
    \item {Economic sustainability}
\end{enumerate}
\vspace{2mm}

\noindent\textbf{Q37. Which business decision provides the most balanced support for the environmental, social, and economic pillars of sustainability?} \hfill [Tag: Pillars of sustainability | Medium]
\begin{enumerate}[label=(\alph*),topsep=2pt,itemsep=1pt,partopsep=0pt,parsep=0pt]
    \item {Expanding production, increasing waste, and funding local advertising}
    \item {Reducing waste, paying fair wages, and maintaining profitability}
    \item {Restoring habitats, ending production, and dismissing all employees}
    \item {Maximizing profits, reducing wages, and meeting minimum emission rules}
\end{enumerate}
\vspace{2mm}

\noindent\textbf{Q38. A program installs solar-powered pumps that provide safe drinking water in villages. Which pair of Sustainable Development Goals is most directly advanced?} \hfill [Tag: Sustainable development goals | Medium]
\begin{enumerate}[label=(\alph*),topsep=2pt,itemsep=1pt,partopsep=0pt,parsep=0pt]
    \item {SDG 6 and SDG 7}
    \item {SDG 4 and SDG 10}
    \item {SDG 9 and SDG 14}
    \item {SDG 12 and SDG 16}
\end{enumerate}
\vspace{2mm}

\noindent\textbf{Q39. A city introduces separate waste collection, repairs reusable goods, and converts organic waste into compost. Which goal is addressed most directly?} \hfill [Tag: Sustainable development goals | Medium]
\begin{enumerate}[label=(\alph*),topsep=2pt,itemsep=1pt,partopsep=0pt,parsep=0pt]
    \item {SDG 8: Decent Work and Economic Growth}
    \item {SDG 14: Life Below Water}
    \item {SDG 16: Peace, Justice and Strong Institutions}
    \item {SDG 12: Responsible Consumption and Production}
\end{enumerate}
\vspace{2mm}

\noindent\textbf{Q40. A government measures progress only by counting newly planted trees, without checking species diversity, survival rates, or restored habitat area. What is the main weakness of this SDG assessment?} \hfill [Tag: Sustainable development goals | Medium]
\begin{enumerate}[label=(\alph*),topsep=2pt,itemsep=1pt,partopsep=0pt,parsep=0pt]
    \item {The indicator measures activity rather than lasting ecological outcomes}
    \item {The indicator focuses more on biodiversity than on tree planting}
    \item {The indicator combines too many independent environmental variables}
    \item {The indicator gives excessive weight to long-term habitat recovery}
\end{enumerate}
\vspace{2mm}

\noindent\textbf{Q41. A coastal city constructs a seawall to reduce storm-surge exposure. The wall protects urban property but accelerates erosion downstream, alters sediment transport, and encourages further construction in the protected zone. Which interpretation best reflects the environmental perspective?} \hfill [Tag: Introduction to environment | Hard]
\begin{enumerate}[label=(\alph*),topsep=2pt,itemsep=1pt,partopsep=0pt,parsep=0pt]
    \item {The wall converts a natural hazard into an interconnected socio-ecological risk}
    \item {The wall reduces total vulnerability whenever protected property values exceed construction costs}
    \item {The wall creates only an economic externality because erosion remains a physical process}
    \item {The wall eliminates environmental risk by separating society from coastal processes}
\end{enumerate}
\vspace{2mm}

\noindent\textbf{Q42. Two valleys receive the same annual rainfall and contain similar soils. One experiences severe flooding after forest conversion and settlement expansion, while the other does not. Which conclusion is most defensible?} \hfill [Tag: Introduction to environment | Hard]
\begin{enumerate}[label=(\alph*),topsep=2pt,itemsep=1pt,partopsep=0pt,parsep=0pt]
    \item {Annual rainfall alone is sufficient to predict the environmental response}
    \item {Environmental outcomes emerge from interactions among natural and human systems}
    \item {Similar soils imply that both valleys must have equal infiltration rates}
    \item {Flooding should be classified as social rather than environmental in the settled valley}
\end{enumerate}
\vspace{2mm}

\noindent\textbf{Q43. A pesticide applied to cropland binds to soil particles, enters a river during erosion, accumulates in fish, and affects fishing-community income. Which sequence of environmental components is most accurate?} \hfill [Tag: Components of environment | Hard]
\begin{enumerate}[label=(\alph*),topsep=2pt,itemsep=1pt,partopsep=0pt,parsep=0pt]
    \item {Social to biotic to abiotic to abiotic}
    \item {Abiotic to biotic to abiotic to social}
    \item {Biotic to abiotic to social to biotic}
    \item {Abiotic to abiotic to biotic to social}
\end{enumerate}
\vspace{2mm}

\noindent\textbf{Q44. An assessment records falling groundwater levels, reduced wetland vegetation, migration of farming households, and a new irrigation-pricing rule. How should these indicators be classified, respectively?} \hfill [Tag: Components of environment | Hard]
\begin{enumerate}[label=(\alph*),topsep=2pt,itemsep=1pt,partopsep=0pt,parsep=0pt]
    \item {Abiotic, biotic, social, institutional}
    \item {Abiotic, social, biotic, institutional}
    \item {Physical, institutional, demographic, biotic}
    \item {Biotic, abiotic, economic, technological}
\end{enumerate}
\vspace{2mm}

\noindent\textbf{Q45. Permafrost thaw releases methane, which strengthens atmospheric warming and causes additional thaw. Which Earth-sphere coupling directly initiates and then amplifies this feedback?} \hfill [Tag: Spheres of earth | Hard]
\begin{enumerate}[label=(\alph*),topsep=2pt,itemsep=1pt,partopsep=0pt,parsep=0pt]
    \item {Hydrosphere-lithosphere coupling followed by lithospheric forcing of the biosphere}
    \item {Atmosphere-biosphere coupling followed by biological forcing of the hydrosphere}
    \item {Cryosphere-atmosphere coupling followed by atmospheric forcing of the cryosphere}
    \item {Biosphere-hydrosphere coupling followed by hydrological forcing of the lithosphere}
\end{enumerate}
\vspace{2mm}

\noindent\textbf{Q46. A volcanic eruption injects sulfur dioxide into the atmosphere, forms sulfate aerosols, reduces incoming solar radiation, and temporarily lowers ocean-surface temperature. Which ordering best represents the dominant cross-sphere pathway?} \hfill [Tag: Spheres of earth | Hard]
\begin{enumerate}[label=(\alph*),topsep=2pt,itemsep=1pt,partopsep=0pt,parsep=0pt]
    \item {Hydrosphere to atmosphere to lithosphere}
    \item {Atmosphere to lithosphere to biosphere}
    \item {Biosphere to hydrosphere to atmosphere}
    \item {Lithosphere to atmosphere to hydrosphere}
\end{enumerate}
\vspace{2mm}

\noindent\textbf{Q47. A proposed dam has positive net monetary benefits but would displace an Indigenous community and fragment a migratory-fish habitat. Which study design best reflects the multidisciplinary nature of environmental studies?} \hfill [Tag: Multidisciplinary nature, scope and importance of environmental studies | Hard]
\begin{enumerate}[label=(\alph*),topsep=2pt,itemsep=1pt,partopsep=0pt,parsep=0pt]
    \item {Integrate ecological connectivity, distributional economics, cultural rights, and hydrology}
    \item {Survey public preferences and replace hydrological measurements with perceived risk}
    \item {Compare construction costs with electricity revenue using a single discount rate}
    \item {Model fish abundance and treat displacement as outside the environmental assessment}
\end{enumerate}
\vspace{2mm}

\noindent\textbf{Q48. Researchers must recommend a groundwater policy under uncertain recharge estimates and unequal access to wells. Which combination of disciplines addresses both the physical limit and the allocation problem most directly?} \hfill [Tag: Multidisciplinary nature, scope and importance of environmental studies | Hard]
\begin{enumerate}[label=(\alph*),topsep=2pt,itemsep=1pt,partopsep=0pt,parsep=0pt]
    \item {Oceanography and genetics supported by architectural history}
    \item {Seismology and epidemiology supported by mineral economics}
    \item {Atmospheric chemistry and taxonomy supported by landscape aesthetics}
    \item {Hydrogeology and environmental economics supported by distributive ethics}
\end{enumerate}
\vspace{2mm}

\noindent\textbf{Q49. A region maintains constant total wealth by replacing a destroyed wetland with financial assets of equal estimated monetary value. The wetland had an irreplaceable flood-regulation threshold. Which sustainability interpretation rejects the substitution?} \hfill [Tag: Concept of sustainability | Hard]
\begin{enumerate}[label=(\alph*),topsep=2pt,itemsep=1pt,partopsep=0pt,parsep=0pt]
    \item {Procedural sustainability because consultation determines ecological equivalence}
    \item {Fiscal sustainability because financial assets cannot fund ecosystem restoration}
    \item {Strong sustainability because critical natural capital must be maintained}
    \item {Weak sustainability because aggregate capital must remain constant}
\end{enumerate}
\vspace{2mm}

\noindent\textbf{Q50. A fishery keeps annual harvest below estimated population growth, yet fish size and genetic diversity decline because only the largest fish are captured. What does this demonstrate?} \hfill [Tag: Concept of sustainability | Hard]
\begin{enumerate}[label=(\alph*),topsep=2pt,itemsep=1pt,partopsep=0pt,parsep=0pt]
    \item {Economic efficiency is sufficient to preserve the fishery's ecosystem functions}
    \item {Genetic change is irrelevant when the total population remains numerically stable}
    \item {A sustainable flow can still degrade the quality and resilience of natural capital}
    \item {Harvest below population growth guarantees ecological sustainability over time}
\end{enumerate}
\vspace{2mm}

\noindent\textbf{Q51. A government subsidizes efficient air conditioners. Electricity used per unit falls by 30\%, but lower operating costs increase ownership and usage enough that total electricity demand rises. Which policy conclusion best follows?} \hfill [Tag: Sustainable development | Hard]
\begin{enumerate}[label=(\alph*),topsep=2pt,itemsep=1pt,partopsep=0pt,parsep=0pt]
    \item {Rising demand is sustainable whenever individual devices consume less electricity}
    \item {The increase proves that efficient technology has no environmental value}
    \item {Efficiency improvements necessarily reduce aggregate resource use over time}
    \item {Efficiency standards should be paired with measures addressing total consumption}
\end{enumerate}
\vspace{2mm}

\noindent\textbf{Q52. Project X raises average income substantially but increases irreversible ecosystem loss and leaves the poorest households worse off. Project Y raises income moderately, preserves ecosystem functions, and improves the position of the poorest. Under an equity-sensitive sustainable development framework, which judgment is strongest?} \hfill [Tag: Sustainable development | Hard]
\begin{enumerate}[label=(\alph*),topsep=2pt,itemsep=1pt,partopsep=0pt,parsep=0pt]
    \item {Project X is preferable because average income dominates non-market considerations}
    \item {Project Y is unsustainable because moderate growth cannot finance future development}
    \item {Project Y is preferable because it better integrates ecological and distributive goals}
    \item {Project X is sustainable if ecosystem losses occur outside the project boundary}
\end{enumerate}
\vspace{2mm}

\noindent\textbf{Q53. An island aquifer receives 12 million m 3 of renewable recharge annually. To prevent saltwater intrusion, withdrawals cannot exceed 75\% of recharge. Domestic demand is 120 m 3 per resident per year, while agriculture requires a fixed 3 million m 3 annually. Ignoring other uses, what is the maximum sustainable population?} \hfill [Tag: Carrying capacity | Hard]
\begin{enumerate}[label=(\alph*),topsep=2pt,itemsep=1pt,partopsep=0pt,parsep=0pt]
    \item {37 ,}
    \item {75 ,}
    \item {100 ,}
    \item {50 ,}
\end{enumerate}
\vspace{2mm}

\noindent\textbf{Q54. A population follows d t d N = r N ( 1 − K N ) with K = 10 , 000 . Habitat loss suddenly reduces K to 6 , 000 while N = 8 , 000 and r > 0 . What does the model predict immediately after the change?} \hfill [Tag: Carrying capacity | Hard]
\begin{enumerate}[label=(\alph*),topsep=2pt,itemsep=1pt,partopsep=0pt,parsep=0pt]
    \item {Negative population growth because the population exceeds the new carrying capacity}
    \item {Zero population growth because carrying capacity affects only long-run equilibrium}
    \item {Unchanged population growth because the intrinsic growth rate has not changed}
    \item {Positive population growth because the population remains below the original capacity}
\end{enumerate}
\vspace{2mm}

\noindent\textbf{Q55. A city imports most of its food and exports its waste, allowing its population to exceed the productive and assimilative capacity of its local land. Which statement best characterizes the city's carrying capacity?} \hfill [Tag: Carrying capacity | Hard]
\begin{enumerate}[label=(\alph*),topsep=2pt,itemsep=1pt,partopsep=0pt,parsep=0pt]
    \item {Its effective capacity equals local food production regardless of imports and waste exports}
    \item {Its local capacity has permanently increased because trade removes ecological constraints}
    \item {Its carrying capacity is undefined because only closed ecosystems possess finite limits}
    \item {Its effective capacity is supported by external ecosystems and may shift impacts elsewhere}
\end{enumerate}
\vspace{2mm}

\noindent\textbf{Q56. A recycling plant is profitable and reduces landfill use, but exposes low-income workers to unsafe pollutant levels. Which sustainability diagnosis is most accurate?} \hfill [Tag: Pillars of sustainability | Hard]
\begin{enumerate}[label=(\alph*),topsep=2pt,itemsep=1pt,partopsep=0pt,parsep=0pt]
    \item {It advances environmental and economic objectives while failing the social pillar}
    \item {It advances the social pillar because recycling creates employment in low-income areas}
    \item {It satisfies all pillars because profitability can compensate workers for environmental risk}
    \item {It fails only the economic pillar because occupational controls increase production costs}
\end{enumerate}
\vspace{2mm}

\noindent\textbf{Q57. A city compares four transport plans. Which plan most clearly represents integration rather than a trade-off among the environmental, social, and economic pillars?} \hfill [Tag: Pillars of sustainability | Hard]
\begin{enumerate}[label=(\alph*),topsep=2pt,itemsep=1pt,partopsep=0pt,parsep=0pt]
    \item {Reliable electric buses with affordable fares funded through congestion pricing}
    \item {A premium metro that reduces emissions but excludes most low-income commuters}
    \item {A low-tax road expansion that reduces costs briefly but induces additional traffic}
    \item {Free diesel buses that improve access but substantially increase local air pollution}
\end{enumerate}
\vspace{2mm}

\noindent\textbf{Q58. A country expands irrigated agriculture to advance SDG 2, but the program depletes aquifers, increases fertilizer runoff, and excludes small farmers. Which SDG principle is most directly violated?} \hfill [Tag: Sustainable development goals | Hard]
\begin{enumerate}[label=(\alph*),topsep=2pt,itemsep=1pt,partopsep=0pt,parsep=0pt]
    \item {National progress is valid whenever one goal improves in absolute terms}
    \item {Goals are integrated and require management of synergies and trade-offs}
    \item {Each goal should be pursued independently through sector-specific optimization}
    \item {Environmental targets take precedence over social targets in every conflict}
\end{enumerate}
\vspace{2mm}

\noindent\textbf{Q59. National access to safely managed drinking water rises from 80\% to 90\% , but access remains 45\% in remote minority communities. Which interpretation best reflects the 2030 Agenda's commitment to leave no one behind?} \hfill [Tag: Sustainable development goals | Hard]
\begin{enumerate}[label=(\alph*),topsep=2pt,itemsep=1pt,partopsep=0pt,parsep=0pt]
    \item {The country should stop collecting subgroup data to preserve universal target reporting}
    \item {The target is effectively achieved because the national average improved by ten points}
    \item {The national gain is insufficient evidence of inclusive progress without disaggregated outcomes}
    \item {The disparity is irrelevant unless national access falls below half the population}
\end{enumerate}
\vspace{2mm}

\noindent\textbf{Q60. A climate project reports avoided emissions but does not count emissions embodied in imported equipment, and it displaces land users without consultation. Which evaluation approach is most consistent with the SDG framework?} \hfill [Tag: Sustainable development goals | Hard]
\begin{enumerate}[label=(\alph*),topsep=2pt,itemsep=1pt,partopsep=0pt,parsep=0pt]
    \item {Evaluate consultation separately and exclude it from sustainable development performance}
    \item {Assess full value-chain effects and social safeguards across interconnected goals}
    \item {Count only territorial emissions because each SDG has an independent accounting boundary}
    \item {Ignore displacement if the project's net carbon balance remains environmentally positive}
\end{enumerate}
\vspace{2mm}


\section{Section B: Comprehensive Subjective Problem with Detailed Analytical Rubric}

\begin{subjectiveq}{1}{1}
\textbf{Ecological Carrying Capacity and IPAT Impact Modelling:}\\A developing urban agglomeration has a population of $P = 3.5 \times 10^6$ residents. The current per-capita economic consumption (affluence) is $A = \text{Rs.~} 250{,}000\text{ per year}$, and the environmental greenhouse gas technology intensity is $T = 0.85\text{ kg } CO_2\text{ equivalent per thousand Rupees of expenditure}$.\\\\(a) Formulate the IPAT equation and compute the current aggregate annual greenhouse gas footprint $I$ in metric megatonnes of $CO_2$ equivalent.\\(b) Demographic and economic projections indicate that over the next 15 years, the metropolitan population will expand by $20\%$ and per-capita affluence will grow by $50\%$. If international and national climate agreements (aligned with SDG 11: Sustainable Cities and SDG 13: Climate Action) mandate a $30\%$ absolute reduction in total emissions $I$ relative to the baseline, calculate the target technological carbon intensity $T_{\text{target}}$ and the required percentage reduction in $T$.\\(c) Provide an analytical discussion on how this technological transition can be engineered through renewable energy integration, transit-oriented development, and circular economy waste frameworks.
\end{subjectiveq}

\begin{solbox}
\textbf{Part (a): Current Annual Carbon Footprint Calculation:}\\The master environmental impact formulation is given by the Ehrlich-Holdren IPAT identity:\begin{equation}I = P \times A \times T\end{equation}Substituting the given baseline values:\begin{align*}P &= 3.5 \times 10^6\text{ persons} \\A &= 250{,}000\text{ Rs./person/year} = 2.5 \times 10^5\text{ Rs./person/year} \\T &= 0.85\text{ kg } CO_2 / 1000\text{ Rs.} = 8.5 \times 10^{-4}\text{ kg } CO_2/\text{Rs.}\end{align*}Evaluating aggregate impact $I$:\begin{align*}I_{\text{baseline}} &= (3.5 \times 10^6) \times (2.5 \times 10^5) \times (8.5 \times 10^{-4})\text{ kg } CO_2/\text{year} \\&= 8.75 \times 10^{11} \times 8.5 \times 10^{-4} = 7.4375 \times 10^8\text{ kg } CO_2/\text{year} \\&= 743{,}750\text{ metric tonnes } CO_2/\text{year} = 0.74375\text{ Megatonnes (Mt) } CO_2/\text{year}.\end{align*}\textbf{Part (b): Computation of Target Technological Intensity $T_{\text{target}}$:}\\Under the projected growth scenarios after 15 years:\begin{align*}P_{\text{new}} &= P \times (1 + 0.20) = 3.5 \times 10^6 \times 1.20 = 4.20 \times 10^6\text{ persons} \\A_{\text{new}} &= A \times (1 + 0.50) = 2.5 \times 10^5 \times 1.50 = 3.75 \times 10^5\text{ Rs./person/year} \\I_{\text{target}} &= I_{\text{baseline}} \times (1 - 0.30) = 0.70 \times 7.4375 \times 10^8 = 5.20625 \times 10^8\text{ kg } CO_2/\text{year}.\end{align*}Setting $I_{\text{target}} = P_{\text{new}} \times A_{\text{new}} \times T_{\text{target}}$:\begin{align*}T_{\text{target}} &= \frac{I_{\text{target}}}{P_{\text{new}} \times A_{\text{new}}} = \frac{5.20625 \times 10^8}{(4.20 \times 10^6) \times (3.75 \times 10^5)} \\&= \frac{5.20625 \times 10^8}{1.575 \times 10^{12}} = 3.30556 \times 10^{-4}\text{ kg } CO_2/\text{Rs.} \\&= 0.33056\text{ kg } CO_2 / 1000\text{ Rs.}\end{align*}The required fractional reduction in technological carbon intensity is:\begin{equation*}\text{Percentage Reduction} = \left(1 - \frac{T_{\text{target}}}{T_{\text{baseline}}}\right) \times 100\% = \left(1 - \frac{3.30556 \times 10^{-4}}{8.5 \times 10^{-4}}\right) \times 100\% = 61.11\%.\end{equation*}\textbf{Part (c): Engineering and Policy Synthesis:}\\To achieve a $61.11\%$ decarbonization efficiency, the municipality must deploy:\begin{itemize}[leftmargin=1.5em]  \item Grid decarbonization: replacing coal-based thermal power with distributed solar PV and battery storage systems.  \item Transport electrification: transitioning bus fleets and two-wheelers to electric mobility backed by metro rail corridors.  \item Energy efficiency building codes: mandatory passive heating/cooling and rooftop green infrastructure.\end{itemize}
\end{solbox}

\begin{examinertip}
\textbf{Marking Scheme Breakdown (Total: 10 Marks):}\\
\textbullet\ \textbf{3 Marks:} Clear formulation of governing environmental identities, equations, and thermodynamic/constitutional foundations.\\
\textbullet\ \textbf{4 Marks:} Rigorous algebraic/numerical substitution, dimensional consistency, and unit handling.\\
\textbullet\ \textbf{3 Marks:} Critical evaluation of environmental implications, regulatory criteria (CPCB/UNDRR/IUCN), and policy mitigations.
\end{examinertip}
